% ============================================================
% РАЗДЕЛ «БЛОК ОТСТРОЙКИ ОТ ПОВРЕЖДЕНИЙ ЗА ТРАНСФОРМАТОРАМИ ОТВЕТВЛЕНИЙ»
% ============================================================

{\color{unimain}\subsection{Блок отстройки от повреждений за трансформаторами ответвлений}}

% Макрос для вставки рисунков с подписью
\newcommand{\optofig}[4][]{%
  \begin{figure}[H]
    \centering
    \includegraphics[#1]{#2}
    \captionof{figure}{#3}\label{#4}
  \end{figure}%
  \medskip
}

\begin{enumerate}

% ============================================================
% ОБЩИЕ СВЕДЕНИЯ
% ============================================================
\item Общие сведения
\begin{enumerate}
\item Для предотвращения ложного срабатывания основной защиты при КЗ за трансформатором ответвления предусмотрен функциональный блок отстройки от повреждений за трансформаторами ответвлений (ОПТО). В качестве измерительных органов ОПТО применяются три междуфазных прямонаправленных реле сопротивления и ОНМНП-Р, блокируемое при выявлении броска намагничивающего тока. При введенном программном ключе контроля защиты от ОПТО подготовка её цепи отключения разрешается только при срабатывании одного из реле сопротивления или ОНМНП-Р, что свидетельствует о наличии КЗ на линии, а не за трансформатором ответвления.
\item Характеристика срабатывания РС (см. рисунок \ref{pict_opto:RS_OPTO}) задается уставками \textbf{«Rср РС-ОПТО»}, \textbf{«Хср РС-ОПТО»} и \textbf{«Ф1 РС-ОПТО»}. Углы наклона характеристики срабатывания органа направленности, обеспечивающего направленность РС, задаются уставками \textbf{«Ф2 ОН-ОПТО»} и \textbf{«Ф3 ОН-ОПТО»}. Для РС в составе ОПТО также предусмотрена возможность выреза в характеристике срабатывания для отстройки от нагрузки. Характеристика срабатывания измерительного органа отстройки от нагрузки определяется уставками \textbf{«Rнг ОПТО»} и \textbf{«Фнг ОПТО»}.
\item Работа ОПТО может быть заблокирована при выявлении неисправности в цепях напряжения при введенном программном ключе \textbf{«БНН ОПТО»}.
\item Вывод ОПТО из работы осуществляется программным переключателем \textbf{«ОПТО»}. При введенном положении программного переключателя \textbf{«ОПТО»} вывод ОПТО из работы осуществляется при наличии сигнала \textit{«Вывод ОПТО»} или \textit{«Вывод терминала»}.
\item Характеристика срабатывания РС в составе ОПТО представлена на рисунке \ref{pict_opto:RS_OPTO}.
  \optofig[width=0.4\textwidth,height=0.4\textheight,keepaspectratio]{\fbpath/02.01 ДЗЛ/561.1491 ОПТО/img/RS_OPTO.pdf}{Характеристика срабатывания РС в составе ОПТО}{pict_opto:RS_OPTO}
\item Функциональная схема отстройки от повреждений за трансформаторами ответвлений представлена на рисунке \ref{pict_opto:OPTO}.
  \optofig[width=1\textwidth,height=1\textheight,keepaspectratio]{\fbpath/02.01 ДЗЛ/561.1491 ОПТО/img/OPTO.pdf}{Функциональная схема отстройки от повреждений за трансформаторами ответвлений}{pict_opto:OPTO}
\item Уставки ОПТО представлены в таблице \ref{app_set:opto}.
\end{enumerate}

\end{enumerate}